% Options for packages loaded elsewhere
\PassOptionsToPackage{unicode}{hyperref}
\PassOptionsToPackage{hyphens}{url}
\documentclass[
  12pt,
]{ctexart}
\usepackage{xcolor}
\usepackage{amsmath,amssymb}
\setcounter{secnumdepth}{-\maxdimen} % remove section numbering
\usepackage{iftex}
\ifPDFTeX
  \usepackage[T1]{fontenc}
  \usepackage[utf8]{inputenc}
  \usepackage{textcomp} % provide euro and other symbols
\else % if luatex or xetex
  \usepackage{unicode-math} % this also loads fontspec
  \defaultfontfeatures{Scale=MatchLowercase}
  \defaultfontfeatures[\rmfamily]{Ligatures=TeX,Scale=1}
\fi
\usepackage{lmodern}
\ifPDFTeX\else
  % xetex/luatex font selection
\fi
% Use upquote if available, for straight quotes in verbatim environments
\IfFileExists{upquote.sty}{\usepackage{upquote}}{}
\IfFileExists{microtype.sty}{% use microtype if available
  \usepackage[]{microtype}
  \UseMicrotypeSet[protrusion]{basicmath} % disable protrusion for tt fonts
}{}
\makeatletter
\@ifundefined{KOMAClassName}{% if non-KOMA class
  \IfFileExists{parskip.sty}{%
    \usepackage{parskip}
  }{% else
    \setlength{\parindent}{0pt}
    \setlength{\parskip}{6pt plus 2pt minus 1pt}}
}{% if KOMA class
  \KOMAoptions{parskip=half}}
\makeatother
\usepackage{longtable,booktabs,array}
\usepackage{calc} % for calculating minipage widths
% Correct order of tables after \paragraph or \subparagraph
\usepackage{etoolbox}
\makeatletter
\patchcmd\longtable{\par}{\if@noskipsec\mbox{}\fi\par}{}{}
\makeatother
% Allow footnotes in longtable head/foot
\IfFileExists{footnotehyper.sty}{\usepackage{footnotehyper}}{\usepackage{footnote}}
\makesavenoteenv{longtable}
\setlength{\emergencystretch}{3em} % prevent overfull lines
\providecommand{\tightlist}{%
  \setlength{\itemsep}{0pt}\setlength{\parskip}{0pt}}
\usepackage{bookmark}
\IfFileExists{xurl.sty}{\usepackage{xurl}}{} % add URL line breaks if available
\urlstyle{same}
\hypersetup{
  hidelinks,
  pdfcreator={LaTeX via pandoc}}

\author{SCX}
\date{2026}

\begin{document}

\section{Theorem 3 (Polished): Unidentifiability of Noise vs.~Learnable
Difficulty}\label{theorem-3-polished-unidentifiability-of-noise-vs.-learnable-difficulty}

\begin{quote}
\textbf{Core claim}: From observable data \((x, y, \{f_m(x)\})\) alone,
label noise and intrinsic sample difficulty are observationally
equivalent. SCX's assumptions (A1)-(A6) are not arbitrary technical
conditions but the minimal sufficient structure needed to break this
equivalence.

\textbf{Revision}: 2026-06-27 --- Generalized to \(K\)-class, arbitrary
state count; added minimal-counterexample proof; explicit
assumption-corollary mapping; positioned within the measurement-error /
label-noise / partial-identification literature.
\end{quote}

\begin{center}\rule{0.5\linewidth}{0.5pt}\end{center}

\subsection{1 Theorem Statement}\label{theorem-statement}

\subsubsection{1.1 Setup}\label{setup}

Let the following be defined as in Theorem 1 (see notation document for
full cross-reference):

\begin{longtable}[]{@{}
  >{\raggedright\arraybackslash}p{(\linewidth - 2\tabcolsep) * \real{0.4706}}
  >{\raggedright\arraybackslash}p{(\linewidth - 2\tabcolsep) * \real{0.5294}}@{}}
\toprule\noalign{}
\begin{minipage}[b]{\linewidth}\raggedright
Symbol
\end{minipage} & \begin{minipage}[b]{\linewidth}\raggedright
Meaning
\end{minipage} \\
\midrule\noalign{}
\endhead
\bottomrule\noalign{}
\endlastfoot
\(\mathcal{X}\) & Input space \\
\(\mathcal{Y}\) & Label space, \(\|\mathcal{Y}\| = K\) \\
\(Y^* \in \mathcal{Y}\) & True label (\textbf{latent}, unobservable) \\
\(Y \in \mathcal{Y}\) & Observed label \\
\(\{f_m\}_{m=1}^M\) & \(M\) expert models \\
\(\mathcal{D}_{\text{obs}}\) & Observable data distribution,
\((X, Y, \{f_m(X)\}) \sim \mathcal{D}_{\text{obs}}\) \\
\(\mathcal{S}\) & State space \\
\(\eta(x) = \mathbb{P}(Y \neq Y^* \mid X = x)\) & Input-dependent noise
rate \\
\end{longtable}

\textbf{What is observable}: The researcher sees only the joint
distribution of \((x_i, y_i, f_1(x_i), \dots, f_M(x_i))\). The true
label \(y^*\), the state \(s(x)\), and the noise indicator are
\textbf{all latent}.

\subsubsection{1.2 Formal Statement}\label{formal-statement}

\textbf{Theorem 3 (Noise-Difficulty Unidentifiability).} For any
\(K \geq 2\) classes, any \(M \geq 1\) experts, and any finite state
space \(\mathcal{S}\), there exist two distinct data-generating
processes \(\mathcal{P}_1\) and \(\mathcal{P}_2\) such that:

\textbf{(i)} Under \(\mathcal{P}_1\), some observed label errors are
caused by \textbf{label noise} (\(y \neq y^*\)).

\textbf{(ii)} Under \(\mathcal{P}_2\), all observed labels equal the
true labels (\(y = y^*\) for all samples), but some samples have
\textbf{intrinsic difficulty}: experts systematically err because of
sample ambiguity rather than label corruption.

\textbf{(iii)} Observational equivalence:

\[\mathcal{P}_1\bigl(x, y, \{f_m(x)\}\bigr) = \mathcal{P}_2\bigl(x, y, \{f_m(x)\}\bigr), \quad \forall (x, y, \{f_m\})\]

\textbf{(iv)} Therefore, no algorithm
\(\mathcal{A}: (\mathcal{X} \times \mathcal{Y} \times \mathcal{Y}^M)^n \to \{0,1\}^n\)
can correctly identify noise in both worlds simultaneously: there exists
\(n\) such that \(\mathcal{A}\) has error rate \(\geq 1/2\) in at least
one world.

\begin{quote}
\textbf{Interpretation}: ``This sample's label is wrong'' and ``This
sample is too hard for the experts'' are observationally equivalent
propositions. SCX's assumption set (A1)-(A6) is precisely what is needed
to break this equivalence.
\end{quote}

\begin{center}\rule{0.5\linewidth}{0.5pt}\end{center}

\subsection{2 Proof: Generalized
Construction}\label{proof-generalized-construction}

\subsubsection{2.1 Design Principles}\label{design-principles}

The construction is designed to be \textbf{minimal}: it uses the fewest
moving parts necessary to establish unidentifiability.

\textbf{Minimality properties}: - \textbf{2 states} suffice (one
ambiguous, one clean): \(\mathcal{S} = \{s_1, s_2\}\) - \textbf{Binary
classification suffices} (\(K = 2\)): the \(K\)-class case is a
corollary - \textbf{1 expert suffices} (\(M = 1\)): the \(M > 1\) case
follows identically - \textbf{No additional structure}: the construction
does not rely on violating any assumption -- it merely shows that within
the space of all possible data-generating processes, indistinguishable
pairs exist

\subsubsection{2.2 World A: Noise-Driven}\label{world-a-noise-driven}

In \(\mathcal{P}_{\text{noise}}\):

\textbf{State \(s_1\)} (``clean and easy''): -
\(\mathbb{P}(X \in s_1) = \rho\), with \(\rho \in (0, 1)\) - Constant
true label: for all \(x \in s_1\), \(y^* = 0\) - Noise mechanism: labels
flipped with probability \(\eta \in (0, 1/2)\)
\[\mathbb{P}(y \neq y^* \mid x \in s_1) = \eta\] - Expert accuracy on
clean samples:
\(\mathbb{P}(f_m(x) = y^* \mid x \in s_1, \text{clean}) = 1 - \varepsilon_1\),
\(\varepsilon_1 \in (0, 1/2)\)

\textbf{State \(s_2\)} (``intrinsically hard''): -
\(\mathbb{P}(X \in s_2) = 1 - \rho\) - Constant true label: for all
\(x \in s_2\), \(y^* = 0\) - No noise:
\(\mathbb{P}(y \neq y^* \mid x \in s_2) = 0\) - Lower expert accuracy:
\(\mathbb{P}(f_m(x) = y^* \mid x \in s_2) = 1 - \varepsilon_2\),
\(\varepsilon_1 < \varepsilon_2 \leq 1/2\)

\subsubsection{2.3 World B:
Difficulty-Driven}\label{world-b-difficulty-driven}

In \(\mathcal{P}_{\text{hard}}\):

\textbf{State \(s_1\)} (``ambiguous'', noise-free): -
\(\mathbb{P}(X \in s_1) = \rho\) (same as World A) - \textbf{All labels
are true}: \(y = y^*\) for every sample - But the true label itself is
ambiguous:
\[\mathbb{P}(y^* = 0 \mid x \in s_1) = 1 - \eta, \quad \mathbb{P}(y^* = 1 \mid x \in s_1) = \eta\]
where \(\eta\) matches World A's noise rate - Expert behavior
(class-conditional): - When \(y^* = 0\):
\(\mathbb{P}(f_m = 0 \mid s_1, y^* = 0) = 1 - \varepsilon_1\) - When
\(y^* = 1\):
\(\mathbb{P}(f_m = 0 \mid s_1, y^* = 1) = 1 - \varepsilon_1\) (biased
toward class 0) - Hence: accuracy is \(1 - \varepsilon_1\) when
\(y^* = 0\), but only \(\varepsilon_1\) when \(y^* = 1\)

\textbf{State \(s_2\)} (``hard'', same as World A): -
\(\mathbb{P}(X \in s_2) = 1 - \rho\), \(y = y^* = 0\) for all samples -
Expert accuracy:
\(\mathbb{P}(f_m(x) = y^* \mid x \in s_2) = 1 - \varepsilon_2\)

\subsubsection{2.4 Observational Equivalence
Verification}\label{observational-equivalence-verification}

Decompose the joint distribution:

\[\mathcal{P}(x, y, f_1, \dots, f_M) = \mathcal{P}(x) \cdot \mathcal{P}(y \mid x) \cdot \prod_{m=1}^M \mathcal{P}(f_m \mid x)\]

\textbf{Step 1: \(\mathcal{P}(x)\)}. \(\mathbb{P}(X \in s_1) = \rho\) in
both worlds. Identical.

\textbf{Step 2: \(\mathcal{P}(y \mid x)\)}.

In \(\mathcal{P}_{\text{noise}}\) (state \(s_1\)): \[\begin{aligned}
\mathcal{P}_{\text{noise}}(y = 0 \mid s_1) &= (1 - \eta) \cdot 1 + \eta \cdot 0 = 1 - \eta \\
\mathcal{P}_{\text{noise}}(y = 1 \mid s_1) &= \eta
\end{aligned}\]

In \(\mathcal{P}_{\text{hard}}\) (state \(s_1\)): \[\begin{aligned}
\mathcal{P}_{\text{hard}}(y = 0 \mid s_1) &= \mathbb{P}(y^* = 0 \mid s_1) = 1 - \eta \\
\mathcal{P}_{\text{hard}}(y = 1 \mid s_1) &= \eta
\end{aligned}\]

State \(s_2\): \(\mathcal{P}(y = 0 \mid s_2) = 1\),
\(\mathcal{P}(y = 1 \mid s_2) = 0\) in both worlds. Identical.
\(\checkmark\)

\textbf{Step 3: \(\mathcal{P}(f_m \mid x)\)}.

In \(\mathcal{P}_{\text{noise}}\) (state \(s_1\)):
\[\mathcal{P}_{\text{noise}}(f_m = 0 \mid s_1) = 1 - \varepsilon_1, \quad \mathcal{P}_{\text{noise}}(f_m = 1 \mid s_1) = \varepsilon_1\]

In \(\mathcal{P}_{\text{hard}}\) (state \(s_1\)): \[\begin{aligned}
\mathcal{P}_{\text{hard}}(f_m = 0 \mid s_1) &= (1-\eta)(1-\varepsilon_1) + \eta(1-\varepsilon_1) = 1 - \varepsilon_1 \\
\mathcal{P}_{\text{hard}}(f_m = 1 \mid s_1) &= (1-\eta)\varepsilon_1 + \eta\varepsilon_1 = \varepsilon_1
\end{aligned}\]

State \(s_2\): \(\mathcal{P}(f_m = 0 \mid s_2) = 1 - \varepsilon_2\),
\(\mathcal{P}(f_m = 1 \mid s_2) = \varepsilon_2\) in both worlds.
Identical. \(\checkmark\)

\textbf{Step 4: Joint identity}. All factors match; therefore:

\[\mathcal{P}_{\text{noise}}(x, y, \{f_m\}) = \mathcal{P}_{\text{hard}}(x, y, \{f_m\}), \quad \forall (x, y, \{f_m\})\]

\textbf{Step 5: Algorithmic indistinguishability}. Any algorithm
\(\mathcal{A}\) receiving \(n\) i.i.d. samples from either distribution
sees identical data. Therefore its output distribution is identical:

\[\mathcal{L}_{\mathcal{P}_{\text{noise}}}(\mathcal{A}(D_n)) = \mathcal{L}_{\mathcal{P}_{\text{hard}}}(\mathcal{A}(D_n))\]

But the true ``noise'' sets differ: - In \(\mathcal{P}_{\text{noise}}\):
\(\mathcal{Z}^*_{\text{noise}} = \{i: x_i \in s_1, y_i = 1\}\) (these
are flipped) - In \(\mathcal{P}_{\text{hard}}\):
\(\mathcal{Z}^*_{\text{hard}} = \varnothing\) (no noise exists)

Since \(\mathcal{A}\) cannot distinguish the worlds, its output must
simultaneously approximate both truth sets. Any deterministic labeling
misclassifies at least half the \(s_1\) samples in at least one world.
Therefore:

\[\max\bigl(\text{Error}_{\mathcal{P}_{\text{noise}}}(\mathcal{A}), \text{Error}_{\mathcal{P}_{\text{hard}}}(\mathcal{A})\bigr) \geq \frac{1}{2}\]

\(\square\)

\subsubsection{2.5 Why This Construction Is
Minimal}\label{why-this-construction-is-minimal}

The construction uses the \textbf{minimum number of parameters} needed
to create observational equivalence between noise and difficulty:

\begin{longtable}[]{@{}
  >{\raggedright\arraybackslash}p{(\linewidth - 4\tabcolsep) * \real{0.2245}}
  >{\raggedright\arraybackslash}p{(\linewidth - 4\tabcolsep) * \real{0.2857}}
  >{\raggedright\arraybackslash}p{(\linewidth - 4\tabcolsep) * \real{0.4898}}@{}}
\toprule\noalign{}
\begin{minipage}[b]{\linewidth}\raggedright
Component
\end{minipage} & \begin{minipage}[b]{\linewidth}\raggedright
Why necessary
\end{minipage} & \begin{minipage}[b]{\linewidth}\raggedright
What happens if removed
\end{minipage} \\
\midrule\noalign{}
\endhead
\bottomrule\noalign{}
\endlastfoot
2 states (\(s_1, s_2\)) & Need one ambiguous state + one clean anchor to
match marginals & With 1 state, the two worlds have different \(P(y)\)
and \(P(f_m)\) marginals, making them distinguishable \\
\(\eta \in (0, 1/2)\) & Creates the proportion of ``suspicious'' samples
& With \(\eta = 0\), both worlds are trivially identical (no noise, no
difficulty) \\
\(\varepsilon_1 \neq \varepsilon_2\) & Creates inter-state expert
accuracy variation & With \(\varepsilon_1 = \varepsilon_2\), the \(s_2\)
anchor no longer constrains the cross-state marginal match \\
Expert bias in World B & Matches the noise-induced expert error pattern
& Without bias, World B's expert marginals would differ from World
A's \\
\end{longtable}

\textbf{No simpler counterexample exists}: Any construction with fewer
states, fewer parameters, or simpler structure cannot simultaneously
match \(\mathcal{P}(y \mid x)\), \(\mathcal{P}(f_m \mid x)\), and
\(\mathcal{P}(x)\) across two causally distinct worlds.

\begin{center}\rule{0.5\linewidth}{0.5pt}\end{center}

\subsection{\texorpdfstring{3 Generalization to \(K\) Classes and
Arbitrary
States}{3 Generalization to K Classes and Arbitrary States}}\label{generalization-to-k-classes-and-arbitrary-states}

\subsubsection{\texorpdfstring{3.1 \(K\)-Class
Construction}{3.1 K-Class Construction}}\label{k-class-construction}

The binary construction extends to general \(K\) as follows:

\textbf{World A (Noise)}: For \(x \in s_1\), true label \(y^* = 0\).
Noise flips labels uniformly over \(\mathcal{Y} \setminus \{0\}\) with
probability \(\eta\):

\[\mathcal{P}_{\text{noise}}(y = 0 \mid s_1) = 1 - \eta, \quad \mathcal{P}_{\text{noise}}(y = c \mid s_1) = \frac{\eta}{K-1}, \; c \neq 0\]

Expert distribution:
\(\mathcal{P}_{\text{noise}}(f_m = 0 \mid s_1) = 1 - \varepsilon_1\),
uniform over other classes.

\textbf{World B (Difficulty)}: True label distribution matches the noise
world's observation profile:

\[\mathcal{P}_{\text{hard}}(y^* = 0 \mid s_1) = 1 - \eta, \quad \mathcal{P}_{\text{hard}}(y^* = c \mid s_1) = \frac{\eta}{K-1}, \; c \neq 0\]

No noise: \(y = y^*\). Expert accuracy: given each \(y^* = c\),
\(\mathbb{P}(f_m = c \mid s_1, y^* = c) = 1 - \varepsilon_1\), errors
uniform over \(K-1\) wrong classes.

\textbf{Verification}: The same marginal-matching argument from Section
2.4 extends directly. The two worlds produce identical
\((X, Y, \{f_m\})\) distributions but have opposite causal
interpretations. \(\square\)

\subsubsection{3.2 Arbitrary State Count}\label{arbitrary-state-count}

The \(K_S\)-state case follows by assigning each state \(s_i\) to one of
two types: - \textbf{Type A} (ambiguous): follows the \(s_1\) pattern
above, with state-specific \(\eta_i\) and \(\varepsilon_{1,i}\) -
\textbf{Type B} (clean anchor): follows the \(s_2\) pattern (no noise,
constant true label)

Observational equivalence holds as long as each Type A state's
parameters are matched across worlds, and Type B states provide
identical constraints. This covers any \(K_S \geq 1\).

\begin{center}\rule{0.5\linewidth}{0.5pt}\end{center}

\subsection{4 Assumption-Identity Mapping: Which A1-A6 Breaks Which
Part}\label{assumption-identity-mapping-which-a1-a6-breaks-which-part}

Each assumption in A1-A6 breaks a specific aspect of Theorem 3's
construction. The mapping below shows that the SCX assumption set is not
arbitrary but \textbf{targeted} at the specific degrees of freedom that
enable unidentifiability.

\subsubsection{4.1 Assumption-to-Construction
Mapping}\label{assumption-to-construction-mapping}

\begin{longtable}[]{@{}
  >{\raggedright\arraybackslash}p{(\linewidth - 4\tabcolsep) * \real{0.1447}}
  >{\raggedright\arraybackslash}p{(\linewidth - 4\tabcolsep) * \real{0.4605}}
  >{\raggedright\arraybackslash}p{(\linewidth - 4\tabcolsep) * \real{0.3947}}@{}}
\toprule\noalign{}
\begin{minipage}[b]{\linewidth}\raggedright
Assumption
\end{minipage} & \begin{minipage}[b]{\linewidth}\raggedright
What It Breaks in the Construction
\end{minipage} & \begin{minipage}[b]{\linewidth}\raggedright
Construction Element Violated
\end{minipage} \\
\midrule\noalign{}
\endhead
\bottomrule\noalign{}
\endlastfoot
\textbf{A1} (Disjoint training) & Expert independence: World B's bias
requires all experts to share the same bias pattern, which is unlikely
if they trained on independent data & Experts' prediction correlation
structure differs between worlds \\
\textbf{A2} (Conditional independence) & The factorization
\(\prod_m \mathcal{P}(f_m \mid x)\) used in verification; without it,
cross-expert correlations provide an extra signal & Joint expert
distribution differs \\
\textbf{A3} (Bounded loss) & Technical (not structural); Hoeffding
inequality requires boundedness & No direct construction violation \\
\textbf{A4} (Uniform independent noise) & World B's difficulty model has
input-dependent label ambiguity (\(\eta\) in \(s_1\) vs \(0\) in
\(s_2\)), violating input-independent noise & Noise rate constancy
across states \\
\textbf{A5} (State homogeneity) & World B's \(s_1\) has two
subpopulations with different expert error rates (\(1-\varepsilon_1\)
for \(y^*=0\) vs \(\varepsilon_1\) for \(y^*=1\)), violating
within-state error uniformity & \(\mu_s\) constancy within state \\
\textbf{A6} (Balanced errors) & World B's experts are biased toward
class 0, concentrating errors on the \(y^*=1 \to \hat{y}=0\) direction,
violating classwise error balance & \(\max_c \mu_c(x) / (\mu_s/(K-1))\)
bounded \\
\end{longtable}

\subsubsection{4.2 Breaking the Construction: Assumption
Sufficiency}\label{breaking-the-construction-assumption-sufficiency}

Each of A1, A4, A5, and A6 alone breaks the specific construction of
Theorem 3. However, this does not mean any single assumption guarantees
global identifiability -- there may exist other unidentifiability
constructions that a single assumption does not break. Hence the need
for the SCX assumption \textbf{set}.

The \textbf{minimal assumption sets} that break all known
unidentifiability constructions are:

\[\begin{aligned}
\text{Set 1: } & \{A1, A4, A5\} \quad \text{(expert independence + uniform noise + state homogeneity)} \\
\text{Set 2: } & \{A1, A4, A6\} \quad \text{(expert independence + uniform noise + balanced errors)} \\
\text{Set 3: } & \{A5, A6, |\mathcal{S}| \geq 2\} \quad \text{(state homogeneity + balanced errors + multi-state structure)}
\end{aligned}\]

SCX uses all six (A1-A6) to ensure robustness across all possible
counterexamples and to enable the quantitative guarantees of Theorem 1.

\subsubsection{4.3 Assumption Violation: What
Degrades}\label{assumption-violation-what-degrades}

\begin{longtable}[]{@{}
  >{\raggedright\arraybackslash}p{(\linewidth - 4\tabcolsep) * \real{0.2703}}
  >{\raggedright\arraybackslash}p{(\linewidth - 4\tabcolsep) * \real{0.4595}}
  >{\raggedright\arraybackslash}p{(\linewidth - 4\tabcolsep) * \real{0.2703}}@{}}
\toprule\noalign{}
\begin{minipage}[b]{\linewidth}\raggedright
Violated
\end{minipage} & \begin{minipage}[b]{\linewidth}\raggedright
Degradation Mode
\end{minipage} & \begin{minipage}[b]{\linewidth}\raggedright
Severity
\end{minipage} \\
\midrule\noalign{}
\endhead
\bottomrule\noalign{}
\endlastfoot
A1 & Expert correlations inflate FPR, Théorème 1's exponential bound
collapses & Severe \\
A4 & State-conditional noise rates bias \(\Delta_s\) estimates &
Moderate \\
A5 & Within-state error heterogeneity makes \(\mu_s\) a poor summary &
Moderate to severe \\
A6 & Concentration of errors inflates FPR in minority classes & Mild to
moderate \\
A3 (bound violation) & Hoeffding/Chernoff inapplicable; need
sub-Gaussian alternatives & Technical \\
All six & Theorem 3's unidentifiability takes full effect &
\textbf{Complete} \\
\end{longtable}

\begin{center}\rule{0.5\linewidth}{0.5pt}\end{center}

\subsection{5 Positioning Within the
Literature}\label{positioning-within-the-literature}

\subsubsection{5.1 Measurement Error Models (Carroll et al.,
2006)}\label{measurement-error-models-carroll-et-al.-2006}

The classical measurement error model posits an unobserved true variable
\(X^*\) and an observed proxy \(X = X^* + U\) where \(U\) is measurement
error. Without validation data or instrumental variables, the joint
distribution of \((X^*, U)\) is \textbf{unidentified}: infinitely many
decompositions of \(P(X)\) into \(P(X^*) * P(U)\) are observationally
equivalent.

Theorem 3 extends this result from the continuous setting to
\textbf{discrete label noise with expert predictions}. The key novelty
is the introduction of \(\{f_m(X)\}\) as auxiliary observations. Theorem
3 shows that even with these extra observations, the identification
problem persists -- the measurement error (noise) and the true label
distribution are not separately identifiable without structural
assumptions.

\textbf{Comparison}:

\begin{longtable}[]{@{}
  >{\raggedright\arraybackslash}p{(\linewidth - 4\tabcolsep) * \real{0.1739}}
  >{\raggedright\arraybackslash}p{(\linewidth - 4\tabcolsep) * \real{0.5870}}
  >{\raggedright\arraybackslash}p{(\linewidth - 4\tabcolsep) * \real{0.2391}}@{}}
\toprule\noalign{}
\begin{minipage}[b]{\linewidth}\raggedright
Aspect
\end{minipage} & \begin{minipage}[b]{\linewidth}\raggedright
Classical Measurement Error
\end{minipage} & \begin{minipage}[b]{\linewidth}\raggedright
Theorem 3
\end{minipage} \\
\midrule\noalign{}
\endhead
\bottomrule\noalign{}
\endlastfoot
Variable type & Continuous \(X^*\) & Discrete \(Y^*\) (labels) \\
Error model & Additive \(X = X^* + U\) & Label flip
\(Y = \text{flip}(Y^*)\) \\
Auxiliary data & Validation data / instruments & Expert predictions
\(\{f_m(X)\}\) \\
Unidentifiability & \(P(X^*) * P(U)\) decomposition & Noise
vs.~difficulty decomposition \\
Resolution & Requires assumptions on \(U\) & Requires A1-A6 (minimum) \\
\end{longtable}

\subsubsection{5.2 Label Noise Theory (Natarajan et al., 2013; Menon et
al., 2015; Patrini et al.,
2017)}\label{label-noise-theory-natarajan-et-al.-2013-menon-et-al.-2015-patrini-et-al.-2017}

The label noise literature establishes that the noise transition matrix
\(T\) and the clean class probabilities \(P(Y^* \mid X)\) are jointly
unidentifiable from \((X, \tilde{Y})\) alone. Standard resolutions
include:

\begin{itemize}
\tightlist
\item
  \textbf{Anchor points} (Liu \& Tao, 2016): points known to belong to a
  specific class with certainty
\item
  \textbf{Noise rate symmetry} (Natarajan et al., 2013): symmetric noise
  with known rate
\item
  \textbf{Dual estimators} (Menon et al., 2015): use of
  class-probability estimators to correct for noise
\end{itemize}

Theorem 3 shows that this unidentifiability \textbf{persists and
deepens} when expert predictions are added. Specifically:

\begin{enumerate}
\def\labelenumi{\arabic{enumi}.}
\tightlist
\item
  Even with observations of \(\{f_m(X)\}\), the noise vs.~difficulty
  ambiguity remains
\item
  The ambiguity is \textbf{semantic}, not just statistical -- two
  fundamentally different causal structures produce the same data
\item
  Resolving this ambiguity requires assumptions about the \textbf{joint
  structure} of experts (A1-A2), not just about the noise mechanism
  itself (A4)
\end{enumerate}

\subsubsection{5.3 Partial Identification (Manski, 2003;
2007)}\label{partial-identification-manski-2003-2007}

Manski's partial identification framework argues that when point
identification is impossible, the researcher should characterize the
\textbf{identified set} -- the set of all parameter values consistent
with the observed distribution.

Theorem 3 can be understood as establishing the identified set for the
noise detection problem: - \textbf{Point estimation} is impossible: a
single ``noise flag'' per sample cannot be consistently estimated - The
\textbf{identified set} contains at least two distinct solutions: the
``noise-driven'' world and the ``difficulty-driven'' world -
\textbf{Assumptions shrink the identified set}: each A1-A6 assumption
reduces the set until point identification becomes possible (Theorem 1
regime) - \textbf{Weak features enlarge the identified set}: even with
A1-A6, weak \(\phi\) makes the SCX estimate indistinguishable from the
loss baseline (Theorem 2 regime)

This positions Theorem 3 within a broader intellectual tradition:
\textbf{identifiability is not binary but graded}, and the SCX framework
explicitly tracks which assumptions are needed to achieve which level of
identification.

\subsubsection{5.4 Graphical Causal Models (Pearl,
2009)}\label{graphical-causal-models-pearl-2009}

From a causal perspective, Theorem 3 identifies a \textbf{Markov
equivalence class} of two causal graphs:

\textbf{Graph A (Noise)}:

\begin{verbatim}
S → Y* → Y ← Noise
↓
X → {f_m}
\end{verbatim}

\textbf{Graph B (Difficulty)}:

\begin{verbatim}
S → Y* (= Y)
↓
X → {f_m}
\end{verbatim}

These graphs share the same skeleton and immoralities (conditional
independence structure) when only \((X, Y, \{f_m\})\) are observed. They
are members of the same Markov equivalence class and cannot be
distinguished by any conditional independence test.

The SCX assumptions break this equivalence by imposing
\textbf{additional constraints} that favor Graph A: - A1 (disjoint
training): restricts the joint expert distribution - A4 (uniform
independent noise): restricts the noise-induced conditional
independencies - A5 (state homogeneity): restricts the within-state
distribution of errors

\subsubsection{5.5 Relationship to Dawid-Skene
Identifiability}\label{relationship-to-dawid-skene-identifiability}

Dawid-Skene (1979) unidentifiability concerns \textbf{label permutation
symmetry}: the EM likelihood has \(K!\) equivalent modes corresponding
to permuting class labels across annotators. This is orthogonal to
Theorem 3's concern.

\begin{longtable}[]{@{}
  >{\raggedright\arraybackslash}p{(\linewidth - 4\tabcolsep) * \real{0.1154}}
  >{\raggedright\arraybackslash}p{(\linewidth - 4\tabcolsep) * \real{0.4615}}
  >{\raggedright\arraybackslash}p{(\linewidth - 4\tabcolsep) * \real{0.4231}}@{}}
\toprule\noalign{}
\begin{minipage}[b]{\linewidth}\raggedright
\end{minipage} & \begin{minipage}[b]{\linewidth}\raggedright
Dawid-Skene
\end{minipage} & \begin{minipage}[b]{\linewidth}\raggedright
Theorem 3
\end{minipage} \\
\midrule\noalign{}
\endhead
\bottomrule\noalign{}
\endlastfoot
\textbf{Ambiguity} & Which permutation of class labels & Noise
vs.~difficulty \\
\textbf{Source} & Permutation symmetry of confusion matrices & Two
causally distinct generative processes \\
\textbf{Resolution} & Anchor annotator, or fix one annotator's confusion
matrix & A1-A6 assumptions \\
\end{longtable}

Even after resolving Dawid-Skene's label permutation problem, Theorem
3's noise-difficulty unidentifiability persists. The two results are
\textbf{orthogonal} and \textbf{cumulative}.

\begin{center}\rule{0.5\linewidth}{0.5pt}\end{center}

\subsection{6 Practical Implications}\label{practical-implications}

\subsubsection{6.1 Justification of the SCX
Framework}\label{justification-of-the-scx-framework}

\begin{enumerate}
\def\labelenumi{\arabic{enumi}.}
\tightlist
\item
  \textbf{A1-A6 are not over-assumptions}: They are the minimal
  structure needed to break a fundamental identifiability failure.
\item
  \textbf{Logical closure}: Theorems 1-3 together cover:

  \begin{itemize}
  \tightlist
  \item
    Necessity (Theorem 3): why assumptions are needed
  \item
    Sufficiency (Theorem 1): what is gained with them
  \item
    Boundary (Theorem 2): when the gain vanishes
  \end{itemize}
\item
  \textbf{Cold-start is theoretically necessary}: Theorem 3 proves that
  initial human review (anchor points) is not a practical compromise but
  a \textbf{theoretical necessity}.
\end{enumerate}

\subsubsection{6.2 What the Practitioner Must
Verify}\label{what-the-practitioner-must-verify}

\begin{longtable}[]{@{}
  >{\raggedright\arraybackslash}p{(\linewidth - 2\tabcolsep) * \real{0.3902}}
  >{\raggedright\arraybackslash}p{(\linewidth - 2\tabcolsep) * \real{0.6098}}@{}}
\toprule\noalign{}
\begin{minipage}[b]{\linewidth}\raggedright
If you claim\ldots{}
\end{minipage} & \begin{minipage}[b]{\linewidth}\raggedright
You must have verified\ldots{}
\end{minipage} \\
\midrule\noalign{}
\endhead
\bottomrule\noalign{}
\endlastfoot
``This sample is label noise'' & A1-A6 hold, or anchor points cover the
relevant state \\
``SCX noise detection works'' & A1-A6 + \(I(\phi; S) \gg 0\) \\
``Noise and difficulty are separated'' & At minimum, A1 + A4 + A5 (or
equivalent) \\
``No assumptions needed'' & Theorem 3 proves this is impossible \\
\end{longtable}

\subsubsection{6.3 Diagnostic Checks for Assumption
Violations}\label{diagnostic-checks-for-assumption-violations}

\begin{longtable}[]{@{}
  >{\raggedright\arraybackslash}p{(\linewidth - 4\tabcolsep) * \real{0.2973}}
  >{\raggedright\arraybackslash}p{(\linewidth - 4\tabcolsep) * \real{0.2973}}
  >{\raggedright\arraybackslash}p{(\linewidth - 4\tabcolsep) * \real{0.4054}}@{}}
\toprule\noalign{}
\begin{minipage}[b]{\linewidth}\raggedright
Assumption
\end{minipage} & \begin{minipage}[b]{\linewidth}\raggedright
Diagnostic
\end{minipage} & \begin{minipage}[b]{\linewidth}\raggedright
Pass threshold
\end{minipage} \\
\midrule\noalign{}
\endhead
\bottomrule\noalign{}
\endlastfoot
A1 (disjoint) & Verify training data provenance & Known disjoint
assignment \\
A4 (uniform noise) & \(\chi^2\) test: \(P(y \neq y^* \mid x \in s)\)
constant across states & \(p > 0.05\) \\
A5 (homogeneity) & KS test: within-state expert error distributions &
\(p > 0.05\) across states \\
A6 (balance) & \(\chi^2\) test: error class distributions balanced &
\(C_{\text{bal}} \leq 2\) \\
\end{longtable}

\subsubsection{6.4 Correct Interpretation of Theorem
3}\label{correct-interpretation-of-theorem-3}

\textbf{What Theorem 3 does NOT say}: - It does not say noise detection
is impossible in practice - It does not say all methods are equally
valid - It does not say SCX's assumptions are the only possible ones

\textbf{What Theorem 3 DOES say}: - Without explicit assumptions, noise
detection lacks a well-defined ground truth - SCX's A1-A6 are exactly
the structure needed - Any method that ``just works'' is implicitly
assuming at least as much

\begin{center}\rule{0.5\linewidth}{0.5pt}\end{center}

\subsection{7 Assumption Necessity
Spectrum}\label{assumption-necessity-spectrum}

The following diagram positions Theorem 3's unidentifiability within a
graded necessity spectrum:

\begin{verbatim}
No assumptions              Partial assumptions            SCX assumptions
    │                              │                              │
    v                              v                              v
Theorem 3              Corollaries 1-6             Theorem 1 guarantees
(full unidentifiability)  (partial identifiability)   (point identification)
    │                              │                              │
    ├── No structure               ├── Anchor points (Cor 1)       ├── A1 + A4 + A5
    │                              ├── Known training (Cor 2)     ├── A1 + A4 + A6
    │                              ├── Uniform noise (Cor 3)     └── A5 + A6 + |S| ≥ 2
    │                              ├── State homogeneity (Cor 4)
    │                              ├── Balanced errors (Cor 5)
    │                              └── Multi-state overlap (Cor 6)
    v                              v                              v
Identified set = all       Identified set = partially      Identified set = single
possible processes         constrained                    point (ground truth)
\end{verbatim}

\begin{center}\rule{0.5\linewidth}{0.5pt}\end{center}

\subsection{References}\label{references}

\subsubsection{Core Theoretical
References}\label{core-theoretical-references}

\begin{enumerate}
\def\labelenumi{\arabic{enumi}.}
\tightlist
\item
  Natarajan, N., Dhillon, I. S., Ravikumar, P. K., \& Tewari, A. (2013).
  Learning with noisy labels. \emph{NeurIPS}.
\item
  Menon, A. K., Van Rooyen, B., Ong, C. S., \& Williamson, R. C. (2015).
  Learning from corrupted binary labels via class-probability
  estimation. \emph{ICML}.
\item
  Patrini, G., Rozza, A., Menon, A. K., Nock, R., \& Qu, L. (2017).
  Making deep neural networks robust to label noise: A loss correction
  approach. \emph{CVPR}.
\end{enumerate}

\subsubsection{Measurement Error Models}\label{measurement-error-models}

\begin{enumerate}
\def\labelenumi{\arabic{enumi}.}
\setcounter{enumi}{3}
\tightlist
\item
  Carroll, R. J., Ruppert, D., Stefanski, L. A., \& Crainiceanu, C. M.
  (2006). \emph{Measurement Error in Nonlinear Models: A Modern
  Perspective} (2nd ed.). Chapman and Hall/CRC.
\item
  Fuller, W. A. (1987). \emph{Measurement Error Models}. Wiley.
\end{enumerate}

\subsubsection{Partial Identification}\label{partial-identification}

\begin{enumerate}
\def\labelenumi{\arabic{enumi}.}
\setcounter{enumi}{5}
\tightlist
\item
  Manski, C. F. (2003). \emph{Partial Identification of Probability
  Distributions}. Springer.
\item
  Manski, C. F. (2007). \emph{Identification for Prediction and
  Decision}. Harvard University Press.
\end{enumerate}

\subsubsection{Graphical Causal Models}\label{graphical-causal-models}

\begin{enumerate}
\def\labelenumi{\arabic{enumi}.}
\setcounter{enumi}{7}
\tightlist
\item
  Pearl, J. (2009). \emph{Causality: Models, Reasoning, and Inference}
  (2nd ed.). Cambridge University Press.
\end{enumerate}

\subsubsection{Finite Mixture Models}\label{finite-mixture-models}

\begin{enumerate}
\def\labelenumi{\arabic{enumi}.}
\setcounter{enumi}{8}
\tightlist
\item
  Teicher, H. (1963). Identifiability of finite mixtures. \emph{The
  Annals of Mathematical Statistics}, 34(4), 1265-1269.
\item
  McLachlan, G. J., \& Peel, D. (2000). \emph{Finite Mixture Models}.
  Wiley.
\end{enumerate}

\subsubsection{Multi-Annotator Models}\label{multi-annotator-models}

\begin{enumerate}
\def\labelenumi{\arabic{enumi}.}
\setcounter{enumi}{10}
\tightlist
\item
  Dawid, A. P., \& Skene, A. M. (1979). Maximum likelihood estimation of
  observer error-rates using the EM algorithm. \emph{JRSS Series C},
  28(1), 20-28.
\item
  Northcutt, C. G., Jiang, L., \& Chuang, I. L. (2021). Confident
  learning: Estimating uncertainty in dataset labels. \emph{JAIR}, 70,
  1373-1411.
\end{enumerate}

\subsubsection{Related SCX Documents}\label{related-scx-documents}

\begin{enumerate}
\def\labelenumi{\arabic{enumi}.}
\setcounter{enumi}{12}
\tightlist
\item
  SCX Theorem 1 (Polished): Multi-Expert Consistency Guarantees.
  \texttt{./01\_noise\_detection\_polished.md}.
\item
  SCX Theorem 2 (Polished): Weak Feature Failure Lower Bound.
  \texttt{./02\_weak\_feature\_polished.md}.
\item
  Unified Notation and Dependencies.
  \texttt{./00\_notation\_and\_dependencies.md}.
\item
  SCX Framework Definitions.
  \texttt{../definitions/01\_state\_conditioned\_risk.md}.
\end{enumerate}

\begin{center}\rule{0.5\linewidth}{0.5pt}\end{center}

\textbf{Revision notes (2026-06-27)}: 1. \textbf{Generalized proof}:
Extracted minimality argument; \(K\)-class construction added; arbitrary
state count addressed. 2. \textbf{Assumption-identity mapping}: New
Section 4 mapping each A1-A6 assumption to the construction element it
breaks. 3. \textbf{Literature connections}: Added Sections 5.1 (Carroll
measurement error), 5.2 (Natarajan label noise), 5.3 (Manski partial
identification), 5.4 (Pearl causality). 4. \textbf{Practical
implications}: New Section 6 with diagnostic checks and
mistaken-interpretation guardrails. 5. \textbf{Assumption necessity
spectrum}: New Section 7 positioning Theorem 3 within a graded
identifiability framework.

\end{document}
